\chapter*{序文}
\addcontentsline{toc}{chapter}{序文}
\chapterepigraph{吾生也有涯，而知也無涯。以有涯隨無涯，殆已。}{『荘子』内篇「養生主」}


ブラックホールに外力を加えると、波の一部は地平面へ吸収され、一部は外部へ
放射される。外部領域だけを系とみなせば、ブラックホールは\term{開いた系}である。
したがって、その摂動を閉じた系の固有振動だけで記述することはできない。本書は、
ブラックホールの線形摂動を、因果的な応答の理論として組み立てる。

本書の基本対象は、外力を応答へ写す\term{因果的な演算子}である。第一章と第二章で
幾何学と場を準備し、第三章でこの演算子を定義する。その後は、同じ演算子を
周波数平面、時間領域、量子論の順に読む。用語を先に並べるのではなく、必要な量を
定義したあとで、その物理的意味を述べる。

この順序により、\term{古典場の散乱}、減衰する\term{共鳴}、\term{量子場の揺らぎ}、孤立量子系の\term{熱化}を
一つの応答理論として接続する。ただし、後の概念を前半の説明へ持ち込まない。
各章は、それ以前に定義した量だけを前提として読むことができるように構成する。

\begin{principlebox}
ブラックホール応答理論の基本対象は、周波数の一覧ではなく、物理的な源を観測される
応答へ写す因果的演算子である。
\end{principlebox}

\section*{想定する読者}

本書は、学部四年生が順に読めることを基準にする。解析力学、電磁気学と波動方程式、
量子力学、熱・統計力学の標準的な学部課程を既習とし、\term{一般相対論}については計量、
共変微分、Schwarzschild 解を見たことがある程度を想定する。\term{ブラックホール摂動論}、
\term{非自己共役スペクトル}、\term{曲がった時空の量子場}、\term{ETH} は前提にしない。

本書は自己完結した数学書を目指さない。\term{特殊関数}と\term{超幾何関数}は、標準的な微分方程式の
解として「初等的な計算」に含めて用いる。一方、\term{解析接続}、\term{積分路変形}、複素
スケーリング、\term{双直交展開}、\term{Riesz 輪郭積分}のように、結論が技法の選び方へ依存するものは、
初出箇所で仕組み、成立条件、数値的な判定基準を説明する。必要な結果を証明しない場合は
仮定と参照先を明記する。この意味での自己整合性を、本書の可読性とする。

本文は参考文献を除いて百頁程度を編集上の目安とするが、厳密な上限にはしない。
定義の順序や因果関係を壊して一頁を削ることはせず、主線から外れる一般論を削ることで
長さを制御する。

\section*{物理ミニマム}

本書の編集原理は、物理的な主線を保つための最小限である。一般相対論、\term{散乱理論}、
\term{非自己共役作用素}、\term{量子場理論}、\term{量子統計}のすべてを網羅することはしない。各概念は、
ブラックホールの応答を一段先へ進める場所でだけ導入する。

本文に残す内容は、少なくとも次の一つを満たさなければならない。

\begin{enumerate}[label=(\roman*),leftmargin=2.8em]
\item 外力と応答の関係を定める。
\item 時間発展を支配する解析構造を定める。
\item 表示に依存する量と観測量を分ける。
\item 古典応答を量子統計の量へ接続する。
\end{enumerate}

この基準を満たさない一般論、長い分類、技法の比較、将来の拡張は、付録または
研究問題へ移す。既知の結果であっても主線に不要ならば削り、未完成の研究を本文の
結論として先取りしない。

\section*{研究問題}

本書では、定型的な練習問題を多数置かない。各部に少数の研究問題を置き、本文で
完成させなかった計算や、これから行うべき研究をそこへ送る。回転ブラックホールの
数値応答、\term{宇宙論的地平面}をもつ系の連続応答、共鳴の\term{高次縮退}、対称性を分解した
熱化の検証は、いずれも一つの研究問題になりうる。

研究問題には、本文から従うことと、まだ保証されていないことを区別して記す。
解答が既知のものについても、完全解答ではなく出発点と判定基準を与えるにとどめる。
問題を解くことは本文の理解の確認ではなく、本文の外へ出ることである。

\section*{本書の構成}

第一部では、亀座標とマスター場から因果的応答を構成する。第四章で準固有振動を
定義し、第五章で合体後のリングダウンを含む時間波形を扱う。ここでブラックホール
摂動論を古典場の応答理論として完成させる。

第二部では、発散する共鳴波を数値的に扱う方法を第六章で定義し、二つの共鳴が
固有ベクトルごと合体する場合を第七章で定義する。第三部では、マスター場を量子化し、
地平面の熱性、熱平衡の二点関数、\term{固有状態熱化仮説}をこの順に導入する。

本書はブラックホールの情報問題、量子カオス、量子情報を網羅する本ではない。
それらが\term{応答関数}に直接触れる場所だけを扱う。この制限によって、古典場から
量子統計までを一つの論理で通し、最後に表示を越えて残る\term{応答不変量}へ戻る。

\section*{Codexとの協働}

本書の構想、章立て、説明順序、章頭引用の選定、初稿の作成、
\LaTeX 組版、および全体の自己整合性の確認には、
OpenAIのAIシステムであるCodexを継続的に用いた。
Codexは単なる校正手段ではなく、構成上および編集上の対話相手として、
本書の形成に参与した。

物理的内容の採否、数式と引用の検証、最終的な記述、および
刊行物に対する責任は、森川億人、小川翔也、廣瀬拓哉が負う。

\vspace{2\baselineskip}
\begin{flushright}
森川億人・小川翔也・廣瀬拓哉\\
2026年8月19日
\end{flushright}

\chapter*{記号を読む前に}
\addcontentsline{toc}{chapter}{記号を読む前に}
\chapterepigraph{舉一隅不以三隅反，則不復也。}{『論語』述而篇 7.8}

本書は式の導入を急ぐ。記号の意味は最初に現れる場所で定め、一般論を先回りして
説明しない。異なる章で同じ記号を使うときは、古典論と量子論で同じ役割を担うように
選ぶ。

序文の章構成は見取り図である。本文では、専門語を原則として定義より前に用いない。
やむを得ず後章の量を予告するときは、定義する章を明記する。

青緑の枠は全体を貫く原理、赤い枠は定義域、解析接続、漸近条件に関する注意である。
灰色の枠は研究問題を表す。研究問題を除き、本文は順に読めば閉じるようにした。
